\chapter*{Executive Summary}
\addcontentsline{toc}{chapter}{Executive Summary}
\label{chap:abstract}
\pagestyle{plain}

An e-commerce application split into four independently deployable microservices
(\texttt{auth-service}, \texttt{products-service}, \texttt{orders-service}, and
\texttt{payments-service}) plus a Next.js frontend, deployed to Oracle Kubernetes
Engine (OKE) through a Terraform-provisioned cluster and a Helm/ArgoCD GitOps
pipeline. Each service maps 1:1 to its own Git repository and container image, and
a push to \texttt{main} reaches a running pod with no manual \texttt{kubectl apply} step.

\begin{itemize}[leftmargin=1.4em]
  \item \textbf{The four services own separate data models but share one database.}
        \texttt{orders-service} and \texttt{payments-service} read and write the
        \texttt{products} collection directly rather than calling
        \texttt{products-service} over HTTP; the backend services do not call each
        other.
  \item \textbf{The infrastructure is fully GitOps-driven.} Each service's CI
        pipeline builds and pushes its image to OCIR, a shared workflow updates the
        matching tag in the Helm chart's \texttt{values.yaml}, and ArgoCD reconciles
        the cluster automatically.
  \item \textbf{Backend mutation routes do not currently enforce authorization}, and
        payment confirmation relies on a client redirect rather than a verified
        Stripe event (see Chapter~\ref{chap:security}).
\end{itemize}

Three storefront pages and the admin dashboard's metrics remain placeholders, and
none of the six repositories currently include automated tests. Chapter~\ref{chap:limitations}
covers the current state in full.
